\chapter{Controle de Fluxo}

\section{\texttt{if} / \texttt{else}}

\begin{lstlisting}[caption={Estrutura básica de if/else}]
let nota = 7.5

if nota >= 7 {
  display "aprovado"
} else if nota >= 5 {
  display "recuperação"
} else {
  display "reprovado"
}
\end{lstlisting}

\subsection*{\texttt{if} como expressão}

\lstinline{if} pode ser usado como expressão — retorna o valor do bloco escolhido:

\begin{lstlisting}
let status = if nota >= 7 { "aprovado" } else { "reprovado" }
display status
\end{lstlisting}

Quando usado como expressão, os dois ramos devem ter o mesmo tipo (ou ambos \texttt{nil}).

\section{\texttt{match}}

\lstinline{match} compara um valor contra múltiplos padrões. O primeiro padrão que
casa é executado:

\begin{lstlisting}[caption={match com literais}]
let codigo = 2

match codigo {
  1 => display "iniciante"
  2 => display "intermediário"
  3 => display "avançado"
  _ => display "desconhecido"
}
\end{lstlisting}

O padrão \lstinline{_} é o caso padrão (\textit{wildcard}) — casa com qualquer valor.

\subsection*{\texttt{match} com strings}

\begin{lstlisting}
let uf = "RS"

match uf {
  "RS" => display "Rio Grande do Sul"
  "SP" => display "São Paulo"
  "RJ" => display "Rio de Janeiro"
  _    => display f"UF desconhecida: {uf}"
}
\end{lstlisting}

\subsection*{\texttt{match} como expressão}

\begin{lstlisting}
let descricao = match codigo {
  1 => "iniciante"
  2 => "intermediário"
  _ => "avançado"
}
\end{lstlisting}

\subsection*{Múltiplos valores por braço}

\begin{lstlisting}
match codigo {
  1 | 2    => display "básico"
  3 | 4    => display "avançado"
  _        => display "outro"
}
\end{lstlisting}

\section{Bloco de expressão}

Um bloco \lstinline|\{ \}| é uma expressão que avalia para o último valor produzido
dentro dele:

\begin{lstlisting}
let resultado = {
  let a = 10
  let b = 20
  a + b      // este é o valor do bloco
}
display resultado    // 30
\end{lstlisting}
